\chapter*{Abstract}
\addcontentsline{toc}{chapter}{Abstract}

Secure Shell (SSH) is a widely used cryptographic protocol for secure remote
administration and data transfer. Its popularity has made it a frequent target
of brute-force attacks, where adversaries attempt repeated login attempts with
different password combinations until access is gained. Such attacks pose
significant risks to organisations, particularly those with limited resources
to deploy advanced intrusion-detection systems.

This project extends a Capstone 1 rule-based detector (static and adaptive
threshold rules applied to Linux authentication logs) with a machine-learning
layer and a hybrid decision fusion engine. The objective is to evaluate
whether a rule + ML hybrid outperforms the pure adaptive-rule baseline on the
same dataset while keeping false positives low and remaining interpretable.

% Methods: real-time log parser + 14 statistical features + 3 classifiers (RF,
% XGBoost, IsolationForest) + a weighted hybrid decision layer.

% Results table is auto-populated by ``scripts/run_pipeline.py`` into
% ``report/figures/baseline_comparison.csv``; drop the table here in Ch.5.

The findings confirm that rule-based detection remains relevant, but its
hybridisation with machine learning further strengthens the practical
defence posture of small organisations and academic institutions.

Keywords: SSH, brute-force detection, authentication logs, anomaly
detection, machine learning.
